% !TeX root = ../教材题库_全部条目_转写工作稿.tex
% Source: XJ-B1-C05-EX-P194-Y1243
% 第5章 / 5.4 函数 y=A sin(ωx+φ) 的图象与性质 / 教材 p.194、195 / PDF 第201、202页
% OCR未校对：公式与文字均可能有误；未完成数学TeX转写；禁止直接发给学生
% 原题及相邻原解截图（核对用）：../题目原图/XJ-B1-C05-EX-P194-Y1243-01.webp, ../题目原图/XJ-B1-C05-EX-P194-Y1243-02.webp
% 本文件只在转写完成后作为题干使用；OCR 条目尚不是可靠题干。
\begin{bankitem}{5.4 例题 例题6}{第5章 / 5.4 函数 y=A sin(\ensuremath{\omega}x+\ensuremath{\varphi}) 的图象与性质 / 教材 p.194、195 / PDF 第201、202页}{XJ-B1-C05-EX-P194-Y1243}{OCR未校对：公式与文字均可能有误；未完成数学TeX转写；禁止直接发给学生}
\begin{ocrdraft}
[教材p.194 / PDF第201页]
\par
0河图5．4一10为某简谐振动的图象，试根
\par
y/cm
\par
据图象回答下列问题：
\par
（1)求该简谐振动的振幅、周期、频率和初相；
\par
3
\par
4
\par
5
\par
6t/s
\par
（2）求t一15s时，振子相对于平衡位置的位移；
\par
（3）写出这个简谐振动的函数解析式．
\par
解（1）从图象上可以看出，这个简谐振动的
\par
图5．4一10
\par
(D在国际单位制中，周期的单位是秒（卜），频率丆的单位是赫兹（Hz），1Hz=1s-l
\par
[教材p.195 / PDF第202页]
\par
振幅为2cm，周期为4，频率为一Hz，初相为0．
\par
（2）由于振动的周期为4s,因此当'=巧s时，位移3，为
\par
一2cm．
\par
（3）设这个简谐振动的函数解析式为
\par
y=Asin()t十，teC0,十亠）．
\par
由一4，得
\par
又A一2，一0
\par
，于是该简谐振动的函数解析式为
\par
3，=2sin
\end{ocrdraft}
\end{bankitem}
